% ═══════════════════════════════════════════════════════════════════════════
% QR-CODES + PLATZKARTEN: Onboarding + Stundenleistung MT-02
% Physik Klasse 8, Doppelstunde 02b
% Enthält: Lernpfad-Onboarding, MT-02, Simulation
% ═══════════════════════════════════════════════════════════════════════════

\documentclass[a4paper,11pt]{article}

\usepackage[utf8]{inputenc}
\usepackage[T1]{fontenc}
\usepackage[ngerman]{babel}
\usepackage{helvet}
\renewcommand{\familydefault}{\sfdefault}

\usepackage[margin=15mm]{geometry}
\usepackage{tikz}
\usepackage{xcolor}
\usepackage{qrcode}
\usepackage{tabularx}
\usepackage{array}

% ═══════════════════════════════════════════════════════════════════════════
% FARBEN
% ═══════════════════════════════════════════════════════════════════════════

\definecolor{physik}{HTML}{2c5aa0}
\definecolor{mittelgrau}{HTML}{888888}
\colorlet{fachfarbe}{physik}

% ═══════════════════════════════════════════════════════════════════════════
% URLs
% ═══════════════════════════════════════════════════════════════════════════

\newcommand{\baseurl}{https://devbox42.github.io/sciencesim/physik/kl08/02-ohm}

% ═══════════════════════════════════════════════════════════════════════════
% DOKUMENT
% ═══════════════════════════════════════════════════════════════════════════

\pagestyle{empty}

\begin{document}

% ───────────────────────────────────────────────────────────────────────────
% SEITE 1: Stundenleistung MT-02 (groß, zum Aufhängen)
% ───────────────────────────────────────────────────────────────────────────

\begin{center}

\vspace*{5mm}
{\fontsize{26}{32}\selectfont\bfseries\color{fachfarbe}Stundenleistung: Ohmsches Gesetz}

\vspace{12mm}

\fcolorbox{fachfarbe}{white}{\qrcode[height=95mm]{\baseurl/MT-02-ohmsches-gesetz.html}}

\vspace{10mm}

{\large\color{gray}devbox42.github.io/.../MT-02-ohmsches-gesetz.html}

\vspace{8mm}

{\Large\color{gray}Scanne den QR-Code mit deinem Gerät}

\vspace{8mm}

\fcolorbox{fachfarbe}{fachfarbe!10}{%
    \parbox{0.8\textwidth}{%
        \centering
        \vspace{3mm}
        \textbf{1.} Wähle deine \textbf{Platznummer} aus\\[2mm]
        \textbf{2.} Wähle dein \textbf{Niveau} ($\star$ / $\star\star$ / $\star\star\star$)\\[2mm]
        \textbf{3.} Bearbeite alle Aufgaben\\[2mm]
        \textbf{4.} Klicke am Ende auf \textbf{„Abgeben"}
        \vspace{3mm}
    }%
}

\vspace{10mm}
{\footnotesize\color{gray}Physik Klasse 8 | Doppelstunde 02b | 35–40 Min}

\end{center}

\newpage

% ───────────────────────────────────────────────────────────────────────────
% SEITE 2: Lernpfad-Onboarding (groß, zum Aufhängen)
% ───────────────────────────────────────────────────────────────────────────

\begin{center}

\vspace*{5mm}
{\fontsize{26}{32}\selectfont\bfseries\color{fachfarbe}Lernpfad: Onboarding Widerstände}

\vspace{12mm}

\fcolorbox{fachfarbe}{white}{\qrcode[height=95mm]{\baseurl/LP-02b-onboarding.html}}

\vspace{10mm}

{\large\color{gray}devbox42.github.io/.../LP-02b-onboarding.html}

\vspace{8mm}

{\Large\color{gray}Scanne den QR-Code mit deinem Gerät}

\vspace{8mm}

\fcolorbox{fachfarbe}{fachfarbe!10}{%
    \parbox{0.75\textwidth}{%
        \centering
        \vspace{3mm}
        Arbeite die 7 Schritte durch und fülle dabei\\[2mm]
        das Arbeitsblatt AB-02b aus.\\[2mm]
        \textit{Vorbereitung für die Stundenleistung!}
        \vspace{3mm}
    }%
}

\vspace{10mm}
{\footnotesize\color{gray}Physik Klasse 8 | Doppelstunde 02b | 30–35 Min}

\end{center}

\newpage

% ───────────────────────────────────────────────────────────────────────────
% SEITE 3: Platzkarten für MT-02 (16 Stück, A4)
% ───────────────────────────────────────────────────────────────────────────

\begin{center}

\vspace*{2mm}
{\fontsize{18}{22}\selectfont\bfseries\color{fachfarbe}Platzkarten – Stundenleistung MT-02}

\vspace{5mm}

% 4×4 Grid mit Platznummern 1-16
\newcommand{\platzkarte}[1]{%
    \fcolorbox{fachfarbe}{white}{%
        \begin{minipage}{42mm}
            \centering
            \vspace{2mm}
            {\footnotesize\color{fachfarbe}Physik Kl. 8}\\[1mm]
            {\Huge\bfseries\color{fachfarbe}#1}\\[2mm]
            \qrcode[height=24mm]{\baseurl/MT-02-ohmsches-gesetz.html}\\[1mm]
            {\tiny MT-02 | Ohmsches Gesetz}
            \vspace{2mm}
        \end{minipage}%
    }%
}

% Zeile 1
\platzkarte{1}\hspace{2mm}\platzkarte{2}\hspace{2mm}\platzkarte{3}\hspace{2mm}\platzkarte{4}

\vspace{3mm}

% Zeile 2
\platzkarte{5}\hspace{2mm}\platzkarte{6}\hspace{2mm}\platzkarte{7}\hspace{2mm}\platzkarte{8}

\vspace{3mm}

% Zeile 3
\platzkarte{9}\hspace{2mm}\platzkarte{10}\hspace{2mm}\platzkarte{11}\hspace{2mm}\platzkarte{12}

\vspace{3mm}

% Zeile 4
\platzkarte{13}\hspace{2mm}\platzkarte{14}\hspace{2mm}\platzkarte{15}\hspace{2mm}\platzkarte{16}

\vspace{8mm}

{\color{gray}\hrule}

\vspace{4mm}

{\footnotesize\color{gray}Ausschneiden und auf die Tische legen. Schüler geben ihre Platznummer im Test ein.}

\end{center}

\newpage

% ───────────────────────────────────────────────────────────────────────────
% SEITE 4: Platzkarten 17-32 (falls benötigt)
% ───────────────────────────────────────────────────────────────────────────

\begin{center}

\vspace*{2mm}
{\fontsize{18}{22}\selectfont\bfseries\color{fachfarbe}Platzkarten – Stundenleistung MT-02 (17–32)}

\vspace{5mm}

% Zeile 1
\platzkarte{17}\hspace{2mm}\platzkarte{18}\hspace{2mm}\platzkarte{19}\hspace{2mm}\platzkarte{20}

\vspace{3mm}

% Zeile 2
\platzkarte{21}\hspace{2mm}\platzkarte{22}\hspace{2mm}\platzkarte{23}\hspace{2mm}\platzkarte{24}

\vspace{3mm}

% Zeile 3
\platzkarte{25}\hspace{2mm}\platzkarte{26}\hspace{2mm}\platzkarte{27}\hspace{2mm}\platzkarte{28}

\vspace{3mm}

% Zeile 4
\platzkarte{29}\hspace{2mm}\platzkarte{30}\hspace{2mm}\platzkarte{31}\hspace{2mm}\platzkarte{32}

\vspace{8mm}

{\color{gray}\hrule}

\vspace{4mm}

{\footnotesize\color{gray}Bei größeren Klassen: Seite 4 zusätzlich drucken.}

\end{center}

\newpage

% ───────────────────────────────────────────────────────────────────────────
% SEITE 5: Kleine QR-Kärtchen (LP + MT + SIM)
% ───────────────────────────────────────────────────────────────────────────

\begin{center}

\vspace*{3mm}
{\fontsize{18}{22}\selectfont\bfseries\color{fachfarbe}QR-Kärtchen zum Ausschneiden}

\vspace{8mm}

% ─── LP-Onboarding ───
\fcolorbox{fachfarbe}{white}{%
    \begin{minipage}{55mm}
        \centering
        \vspace{3mm}
        {\small\bfseries\color{fachfarbe}Lernpfad Onboarding}\\[3mm]
        \qrcode[height=38mm]{\baseurl/LP-02b-onboarding.html}\\[2mm]
        {\tiny LP-02b-onboarding | 7 Schritte}
        \vspace{3mm}
    \end{minipage}%
}
\hspace{8mm}
% ─── MT-02 ───
\fcolorbox{fachfarbe}{white}{%
    \begin{minipage}{55mm}
        \centering
        \vspace{3mm}
        {\small\bfseries\color{fachfarbe}Stundenleistung MT-02}\\[3mm]
        \qrcode[height=38mm]{\baseurl/MT-02-ohmsches-gesetz.html}\\[2mm]
        {\tiny MT-02 | Ohmsches Gesetz}
        \vspace{3mm}
    \end{minipage}%
}
\hspace{8mm}
% ─── SIM-02b ───
\fcolorbox{fachfarbe}{white}{%
    \begin{minipage}{55mm}
        \centering
        \vspace{3mm}
        {\small\bfseries\color{fachfarbe}Simulation Reihenschaltung}\\[3mm]
        \qrcode[height=38mm]{\baseurl/SIM-02b-widerstand-reihe.html}\\[2mm]
        {\tiny SIM-02b | R in Reihe}
        \vspace{3mm}
    \end{minipage}%
}

\vspace{15mm}

{\color{gray}\hrule}

\vspace{6mm}

% Nochmals 3er-Set zum Ausschneiden
{\small\color{fachfarbe}\bfseries Zusätzliche Kärtchen (gleicher Inhalt)}

\vspace{6mm}

\fcolorbox{fachfarbe}{white}{%
    \begin{minipage}{55mm}
        \centering
        \vspace{3mm}
        {\small\bfseries\color{fachfarbe}Lernpfad Onboarding}\\[3mm]
        \qrcode[height=38mm]{\baseurl/LP-02b-onboarding.html}\\[2mm]
        {\tiny LP-02b-onboarding | 7 Schritte}
        \vspace{3mm}
    \end{minipage}%
}
\hspace{8mm}
\fcolorbox{fachfarbe}{white}{%
    \begin{minipage}{55mm}
        \centering
        \vspace{3mm}
        {\small\bfseries\color{fachfarbe}Stundenleistung MT-02}\\[3mm]
        \qrcode[height=38mm]{\baseurl/MT-02-ohmsches-gesetz.html}\\[2mm]
        {\tiny MT-02 | Ohmsches Gesetz}
        \vspace{3mm}
    \end{minipage}%
}
\hspace{8mm}
\fcolorbox{fachfarbe}{white}{%
    \begin{minipage}{55mm}
        \centering
        \vspace{3mm}
        {\small\bfseries\color{fachfarbe}Simulation Reihenschaltung}\\[3mm]
        \qrcode[height=38mm]{\baseurl/SIM-02b-widerstand-reihe.html}\\[2mm]
        {\tiny SIM-02b | R in Reihe}
        \vspace{3mm}
    \end{minipage}%
}

\vspace{10mm}

{\footnotesize\color{gray}Seite 1+2: groß aufhängen | Seite 3+4: Platzkarten | Seite 5: für Gruppentische}

\end{center}

\end{document}
